\begin{figure}[t]
\centering
\begin{tikzpicture}
\begin{axis}[
    width=\columnwidth,
    height=7.5cm,
    xlabel={$\beta$ (zi$^{-1}$)},
    ylabel={$\gamma$ (zi$^{-1}$)},
    xmin=0.05, xmax=0.6,
    ymin=0.05, ymax=0.40,
    colormap/viridis,
    colorbar,
    colorbar style={
        title={$\log_{10}\mathrm{MSE}$},
        title style={font=\footnotesize},
        ticklabel style={font=\footnotesize},
    },
    point meta min=6.9,
    point meta max=10,
    tick label style={font=\footnotesize},
    label style={font=\small},
]
% peisajul de pierdere
\addplot[
    scatter, only marks, mark=square*, mark size=2.7pt,
    scatter src=explicit,
] table[col sep=comma, x=beta, y=gamma, meta=log_mse]
    {data/identifiability_grid_sir.csv};

% optimul liber
\addplot[mark=*, mark size=4pt, color=red, mark options={fill=white, line width=1pt}]
    coordinates {(0.3677, 0.3346)};
\node[font=\scriptsize, anchor=south, color=red, fill=white, fill opacity=0.85, text opacity=1, inner sep=1pt]
    at (axis cs:0.3677, 0.3346) {liber};

% optimul constrans (gamma=1/14)
\addplot[mark=*, mark size=4pt, color=red!70!black, mark options={fill=white, line width=1pt}]
    coordinates {(0.1039, 0.0714)};
\node[font=\scriptsize, anchor=south, color=red!70!black, fill=white, fill opacity=0.85, text opacity=1, inner sep=1pt]
    at (axis cs:0.1039, 0.0714) {constrans};

% linia gamma=1/14
\draw[dashed, red!50, thick] (axis cs:0.05, 0.0714) -- (axis cs:0.6, 0.0714);
\node[font=\scriptsize, anchor=west, color=red!70!black]
    at (axis cs:0.55, 0.0714) {$\gamma = 1/14$};

\end{axis}
\end{tikzpicture}
\caption{Peisajul de pierdere $\log_{10}\mathrm{MSE}(\beta, \gamma)$ pentru
modelul SIR pe fereastra de antrenare. Punctele rosii marcheaza solutia
libera $(0{,}368, 0{,}335)$ si solutia constransa $(0{,}104, 0{,}0714)$;
linia rosie punctata indica restrictia clinica $\gamma = 1/14$. Valea de
pierdere mica este alungita si aproape liniara, ceea ce inseamna ca multe
combinatii $(\beta, \gamma)$ produc o potrivire similara cu datele observate
din valul 1 --- semnatura tipica a identifiabilitatii practice slabe pentru
SIR cu o singura serie observata.}
\label{fig:identifiability_landscape}
\end{figure}
